% Cayley's Collected Mathematical Papers, Vol. 10, Pages 101--150
% Transcribed from: collectedmathema10cayluoft.pdf
% Papers 654 (conclusion) and 655

\documentclass[12pt,a4paper]{article}

% ==================== STANDARD ENGLISH MATH PREAMBLE ====================
\usepackage[utf8]{inputenc}
\usepackage[T1]{fontenc}
\usepackage[english]{babel}
\usepackage{amsmath,amssymb,amsthm}
\usepackage{amsfonts}
\usepackage{mathtools}
\usepackage{bm}
\usepackage{array}
\usepackage{booktabs}
\usepackage{geometry}
\usepackage{fancyhdr}
\usepackage{titlesec}
\usepackage{enumitem}
\usepackage{mathrsfs}
\usepackage{graphicx}

\geometry{margin=1in}

% Theorem-like environments
\theoremstyle{plain}
\newtheorem{theorem}{Theorem}
\newtheorem{lemma}{Lemma}
\newtheorem{corollary}{Corollary}

% Custom commands for frequently used notation
\newcommand{\dd}{\mathrm{d}}
\newcommand{\pd}[2]{\frac{\partial #1}{\partial #2}}
\newcommand{\pdd}[2]{\frac{\partial^{2} #1}{\partial #2^{2}}}
\DeclareMathOperator{\Jac}{Jac}

% Header/Footer
\pagestyle{fancy}
\fancyhf{}
\fancyhead[C]{\small Cayley's Collected Mathematical Papers --- Vol.~10}
\fancyfoot[C]{\thepage}
\renewcommand{\headrulewidth}{0.4pt}

% Section formatting
\titleformat{\section}{\large\bfseries}{\thesection.}{0.5em}{}
\titleformat{\subsection}{\normalsize\bfseries}{\thesubsection.}{0.5em}{}

\begin{document}

% =======================================================================
% TITLE PAGE INFORMATION
% =======================================================================
\begin{center}
{\Large\bfseries The Collected Mathematical Papers of\\[0.3em] Arthur Cayley, F.R.S.}\\[1em]
{\large Volume 10}\\[2em]
{\large Pages 101--150}\\[1em]
\rule{0.6\textwidth}{0.4pt}\\[1em]
{\normalsize Paper 654 (conclusion): \textit{On Certain Octic Surfaces}}\\[0.3em]
{\normalsize Paper 655: \textit{A Memoir on Differential Equations}}\\[2em]
\end{center}

\thispagestyle{empty}
\newpage
\setcounter{page}{101}

% #############################################################################
% PAPER 654: ON CERTAIN OCTIC SURFACES (conclusion)
% From Quarterly Journal of Pure and Applied Mathematics, vol. VII (1867), pp. 164--174
% Original article pages 81--92
% #############################################################################
\section*{654. On Certain Octic Surfaces.}
\addcontentsline{toc}{section}{654. On Certain Octic Surfaces}

\begin{center}
[From the \textit{Quarterly Journal of Pure and Applied Mathematics,} vol.~\textsc{vii} (1867), pp.~164--174.]
\end{center}

\subsection*{Conclusion of Paper}

\setcounter{equation}{0}
\setcounter{page}{101}

In fact, the equation of the surface may be written in the form
%
\begin{multline}
w^{2} \{f x^{2} + g y^{2} + h z^{2} - 2gh y^{2} z^{2} - 2hf z^{2} x^{2} - 2fg x^{2} y^{2}\} \\
+ 2w \{-cf x^{2} - ag y^{2} - bh z^{2} x^{2} + 2k a^{2} b^{2} c^{2} x^{2} y^{2} z^{2}
+ bf x^{2} z^{2} + cg y^{2} x^{2} + ah z^{2} y^{2}\} \\
+ \{ay^{2} z^{2} + bz^{2} x^{2} + cx^{2} y^{2}\}^{2} = 0,
\end{multline}
%
which puts in evidence the nodal curve
\[
w = 0, \quad -ay^{2} z^{2} - bz^{2} x^{2} - cx^{2} y^{2} = 0:
\]
there are three similar forms which put in evidence the other three nodal curves.

The four curves are so related to each other that every line which meets three of them meets also the fourth curve; there is consequently a singly infinite series of lines meeting each of the four curves; these break up into four series of lines each forming an octic scroll, and each scroll has the four curves for nodal curves respectively; that is, each scroll is a surface included under the foregoing general equation, and derived from it by assigning a proper value to the constant~$k$.

To determine these values, write
\[
\frac{\lambda}{af} + \frac{\mu}{bg} + \frac{\nu}{ch} = 0, \qquad
\frac{\lambda^{2}}{af} + \frac{\mu^{2}}{bg} + \frac{\nu^{2}}{ch} = 0,
\]
equations which give four systems of values for the ratios $(\lambda : \mu : \nu)$. We have then
\[
k = af \frac{\mu\,\nu}{\lambda^{2}} + bg \frac{\nu\,\lambda}{\mu^{2}} + ch \frac{\lambda\,\mu}{\nu^{2}},
\]
viz.~$k$ has four values corresponding to the several values of $(\lambda : \mu : \nu)$.

The scroll in question is M.~De~La~Gourn\'{e}rie's scroll $\Sigma_{1}$; the equation of the scroll $\Sigma_{1}$ is consequently obtained from the octic equation by writing therein the last-mentioned value of~$k$.

It is to be noticed that $k$ is, in effect, determined by a quartic equation; and, that, for a certain relation between the coefficients, this equation will have a twofold root. Assuming that this relation is satisfied, and assigning to $k$ its twofold value, the resulting scroll becomes a torse; that is, two of the four scrolls coincide together and degenerate into a torse; corresponding to the remaining two values of $k$ we have two scrolls, companions of the torse. In order to a twofold value of $k$, we must have
\[
\frac{af}{\lambda} = \frac{bg}{\mu} = \frac{ch}{\nu},
\]
and thence
\[
(af)^{\frac{1}{3}} + (bg)^{\frac{1}{3}} + (ch)^{\frac{1}{3}} = 0;
\]
or, what is the same thing,
\[
(af + bg + ch)^{3} - 27abcfgh = 0.
\]

\medskip\noindent\textsc{C.~X.}\hfill 11

\newpage
\setcounter{page}{102}

If for a moment we write
\[
af = \alpha^{3}, \quad bg = \beta^{3}, \quad ch = \gamma^{3}, \quad \text{and therefore, } \alpha + \beta + \gamma = 0;
\]
then for the twofold root, we have $\lambda : \mu : \nu = \alpha : \beta : \gamma$, and consequently
\begin{align*}
k &= \alpha^{3}(\gamma - \beta) + \beta^{3}(\alpha - \gamma) + \gamma^{3}(\beta - \alpha) \\
&= (\alpha - \beta)(\beta - \gamma)(\gamma - \alpha),
\end{align*}
that is,
\[
k = \{(af)^{\frac{1}{3}} - (bg)^{\frac{1}{3}}\}\,\{(bg)^{\frac{1}{3}} - (ch)^{\frac{1}{3}}\}\,\{(ch)^{\frac{1}{3}} - (af)^{\frac{1}{3}}\},
\]
which agrees with the result in regard to the octic torse.

If in the octic equation we write $(x, y, z, w)$ in place of $(x^{2}, y^{2}, z^{2}, w^{2})$, then we have the quartic equation
%
\begin{multline}
fx^{2} + gy^{2} + hz^{2} + 2bc yz + 2ca zx + 2ab xy \\
- 2fgh y^{2} z^{2} - 2hfa z^{2} x^{2} - 2fgb x^{2} y^{2} \\
+ 2cay^{2}zx - 2agy^{2}zw + 2cgz^{2}xw \\
+ 2abz^{2}xy - 2bhz^{2}xw + 2ahz^{2}yw \\
- 2ghw^{2}yz - 2hfw^{2}zx - 2fgw^{2}xy \\
+ 2kxyzw = 0,
\end{multline}
%
which is the equation of a quartic surface touched by the planes $x=0$, $y=0$, $z=0$, $w=0$, in the four conics
\begin{align*}
x&=0, & hzw - gwy + ayz &= 0, \\
y&=0, & -hzw + fwx + bzx &= 0, \\
z&=0, & gyw - fwx + cxy &= 0, \\
w&=0, & -ayz - bzx - cxy &= 0,
\end{align*}
respectively.

\subsection*{II.}

The octic surface
%
\begin{multline}
U = b^{2}c^{2}f^{2}x^{4} + c^{2}a^{2}g^{2}y^{4} + a^{2}b^{2}h^{2}z^{4} + f^{2}g^{2}h^{2}w^{4} \\
- 2a^{2}cg(bg - ch)f^{2}y^{2}z^{2} - 2b^{2}ah(ch - af)g^{2}z^{2}x^{2} - 2c^{2}bf(af - bg)h^{2}x^{2}y^{2} \\
+ 2a^{2}bh(\,\,)y^{2}w^{2} + 2b^{2}cf(\,\,)z^{2}w^{2} + 2c^{2}ag(\,\,)x^{2}w^{2} \\
- 2f^{2}bc(\,\,)x^{2}v^{2} - 2g^{2}ca(\,\,)y^{2}v^{2} - 2h^{2}ab(\,\,)z^{2}v^{2} \\
+ 2fghi(\,\,)x^{2}w^{2} + 2g^{2}hf(\,\,)y^{2}v^{2} + 2h^{2}fg(\,\,)z^{2}w^{2} \\
+ f^{2}(bg + ch - 4bgch)w^{2}x^{2} + g^{2}(ch + af - 4chaf)w^{2}y^{2} + h^{2}(af + bg - 4ahfg)w^{2}z^{2} \\
+ a^{2}(\,\,)y^{2}z^{2} + b^{2}(\,\,)z^{2}x^{2} + c^{2}(\,\,)x^{2}y^{2} \\
- 2gh(bcgh - a^{2}f^{2} - 2afbg - 2afch)w^{2}y^{2}z^{2} \\
- 2bh(\,\,)z^{2}x^{2}w^{2} + 2cg(\,\,)y^{2}x^{2}w^{2} + 2bc(\,\,)x^{2}y^{2}z^{2} \\
- 2hf(cahf - bg - 2bgaf - 2bgch)w^{2}z^{2}x^{2} \\
- 2cf(\,\,)x^{2}y^{2}w^{2} + 2ah(\,\,)z^{2}y^{2}w^{2} + 2ca(\,\,)y^{2}x^{2}z^{2} \\
- 2fg(abfg - ch - 2chaf - 2chbg)w^{2}x^{2}y^{2} \\
- 2ag(\,\,)y^{2}z^{2}w^{2} + 2bf(\,\,)x^{2}z^{2}w^{2} + 2ab(\,\,)z^{2}x^{2}y^{2} \\
+ 2\Omega x^{2}y^{2}z^{2}w^{2} = 0,
\end{multline}
%
where the values of the coefficients indicated by $(\,\,)$ are at once obtained by the proper interchanges of the letters, and where $\Omega$ is an arbitrary coefficient, is a surface having the four nodal conics
\begin{align*}
x&=0, & cf y^{2} - bz^{2} + fw^{2} &= 0, \\
y&=0, & -ca z^{2} + az^{2} + gw^{2} &= 0, \\
z&=0, & ba y^{2} - ay^{2} + hw^{2} &= 0, \\
w&=0, & -fax^{2} - gy^{2} - hz^{2} &= 0.
\end{align*}

\newpage
\setcounter{page}{103}

In fact, writing the equation under the form
\[
w^{2}\Theta + (fax^{2} + gy^{2} + hz^{2})^{2} \times (b^{2}c^{2}y^{2}z^{2} + c^{2}a^{2}z^{2}x^{2} + a^{2}b^{2}x^{2}y^{2} - 2a^{2}bc y^{2}z^{2} - 2b^{2}ca z^{2}x^{2} - 2c^{2}ab x^{2}y^{2}) = 0,
\]
we put in evidence the nodal conic $w = 0$, $fax^{2} + gy^{2} + hz^{2} = 0$: and similarly for the other nodal conics.

It is to be observed, that the complete section by the plane $w = 0$ is the conic $fax^{2} + gy^{2} + hz^{2} = 0$, twice repeated, and the quartic
\[
b^{2}c^{2}x^{4} + c^{2}a^{2}y^{4} + a^{2}b^{2}z^{4} - 2a^{2}bc y^{2}z^{2} - 2b^{2}ca z^{2}x^{2} - 2c^{2}ab x^{2}y^{2} = 0:
\]
the latter being the system of four lines
\[
\frac{x}{a} \pm \frac{y}{b} \pm \frac{z}{c} = 0.
\]
The plane in question, $w = 0$, meets the other nodal conics in the six points
\[
(x=0,\; -b^{2} + cz^{2}=0), \quad (y=0,\; cz^{2} - ax^{2}=0), \quad (z=0,\; cx^{2} - by^{2}=0),
\]
which six points are the angles of the quadrilateral formed by the above-mentioned four lines.

The four conics are such, that every line meeting three of these conics meets also the fourth conic. The lines in question form a double system: each of these

\medskip\noindent 11---2

\newpage
\setcounter{page}{104}

systems has, in reference to any pair of nodal conics, a homographic property as follows; viz. considering for example the two conics in the planes $z = 0$ and $w = 0$ respectively, if a line meets these conics in the points $P$ and $Q$ respectively, and through these points respectively and the line $x=0$, $y=0$ we draw planes, then the system of the $P$~planes and the system of the $Q$~planes correspond homographically to each other, the coincident planes of the two systems being the planes $x=0$ and $y=0$ respectively.

Conversely, if through the line $(x=0,\; y=0)$ we draw the two homographically related planes meeting the two conics in the points $P$ and $Q$ respectively, then, for a proper value (determined by a quadratic equation) of the constant $k(=\theta' \div \theta)$ which determines the homographic relation, the line $PQ$ will be a line meeting each of the four conics, and will belong to one or other of the above-mentioned two systems, as $k$ is equal to one or the other of the two roots of the quadratic equation.

The scroll generated by the lines meeting each of the four conics, or what is the same thing, any three of these conics, is prim\^{a} facie a scroll of the order~16; but by what precedes, it appears that this scroll breaks up into two scrolls, which will be each of the order~8. Moreover, each scroll has the four conies for nodal curves; and since the equation $U=0$ is the general equation of an octic surface having the four conics for nodal curves, it follows, that the equation of the scroll is derived from that of the octic surface $U=0$, by assigning a proper value to the indeterminate coefficient $\Omega$; so that there are in fact two values of $\Omega$, for each of which the surface $U=0$ becomes a scroll.

To sustain the foregoing conclusions, take $x = \theta'y$, $x = \theta y$ for the equations of the two planes through the line $(x=0,\; y=0)$, which meet the $z$-conic and $w$-conic in the points $P$ and $Q$ respectively; then the equations of the line $PQ$ are
\begin{align*}
\sqrt{}(h\theta^{2} + g)(x - \theta'y) + \sqrt{}(-h)(\theta' - \theta)z &= 0, \\
-\sqrt{}(h\theta'^{2} - a)(x - \theta y) + \sqrt{}(-h)(\theta' - \theta)w &= 0,
\end{align*}
or, writing therein $\theta' = k\theta$, the equations are
\begin{align*}
\sqrt{}(h\theta^{2} + g)(x - k\theta y) + \sqrt{}(-h)(k - 1)\theta z &= 0, \\
-\sqrt{}(hk^{2}\theta^{2} - a)(x - \theta y) + \sqrt{}(-h)(k - 1)\theta w &= 0.
\end{align*}
To find where the line in question meets the plane $y = 0$, we have
\begin{align*}
\sqrt{}(h\theta^{2} + g)x + \sqrt{}(-h)(k - 1)\theta z &= 0, \\
-\sqrt{}(hk^{2}\theta^{2} - a)x + \sqrt{}(-h)(k - 1)\theta w &= 0,
\end{align*}
and thence
\begin{align*}
(h\theta^{2} + g)x^{2} + h(k-1)^{2}\theta^{2}z^{2} &= 0, \\
(hk^{2}\theta^{2} - a)x^{2} + h(k-1)^{2}\theta^{2}w^{2} &= 0,
\end{align*}
or multiplying $a$, $g$ and adding
\[
(af + hgk^{2})x^{2} + h(k-1)^{2}(az^{2} + gw^{2}) = 0,
\]
or assuming
\[
af + hgk^{2} + ch(k-1)^{2} = 0,
\]
the equation is
\[
-cx^{2} + az^{2} + gw^{2} = 0.
\]

\newpage
\setcounter{page}{105}

That is, $k$ being determined by the quadric equation
\[
af + hgk^{2} + ch(k-1)^{2} = 0,
\]
the line $PQ$ meets the $y$-conic $y=0$, $-ca^{2} + az^{2} + gw^{2} = 0$; and, in a similar manner, it appears that the line $PQ$ also meets the $x$-conic $x=0$, $cy^{2} - hz^{2} + fw^{2} = 0$.

Writing for greater symmetry
\[
\lambda : -k : k-1 = \lambda : \mu : \nu,
\]
we have
\[
\lambda + \mu + \nu = 0, \qquad af\lambda^{2} + hg\mu^{2} + ch\nu^{2} = 0,
\]
so that there are two systems of values of $(\lambda, \mu, \nu)$ corresponding to, and which may be used in place of, the two values of~$k$ respectively.

Starting now from the equations
\begin{align*}
(h\theta^{2} + g)(k\theta y - x)^{2} + h(k-1)^{2}\theta^{2}z^{2} &= 0, \\
(bk^{2}\theta^{2} - a)(\theta y - x)^{2} + h(k-1)^{2}\theta^{2}w^{2} &= 0,
\end{align*}
the elimination of $\theta$ from these equations leads to an equation $U=0$, of the above mentioned form but with a determinate value of the coefficient.

The process, although a long one, is interesting and I give it in some detail.

\subsubsection*{Elimination of $\theta$ from the foregoing equations.}

We have
\[
U = \text{M}\,\Pi\,[(h\theta^{2} + g)(k\theta y - x)^{2} + h(k-1)^{2}\theta^{2}z^{2}],
\]
where $\Pi$ denotes the product of the expressions corresponding to the four roots of the equation
\[
(bk^{2}\theta^{2} - a)(\theta y - x)^{2} + h(k-1)^{2}\theta^{2}w^{2} = 0.
\]
Observing that this equation does not contain $z$, and that the expression under the sign $\Pi$ does not contain $w$, it is at once seen that the product $\Pi$ is in regard to $(z, w)$ a rational and integral function of the form $(z^{2}, w^{2})^{2}$; and since, in regard to $(z, w)$, $U$ is also a rational and integral function of the same form $(z^{2}, w^{2})^{2}$, it is clear that the factor M does not contain $z$ or $w$, but is a function of only $(x, y)$.

To determine it we may write $z=0$, $w=0$: this gives
\[
c^{2}(bk^{2}\theta^{2} - a)^{2}y^{2}(fx^{2} + gy^{2}) = \text{M}\,\Pi\,(f\theta^{2} + g)(k\theta y - x)^{2},
\]
where
\[
(bk^{2}\theta^{2} - a)(\theta y - x)^{2} = 0,
\]
and the values of $\theta$ are therefore
\[
+\sqrt{\frac{a}{bk}}\,\frac{x}{y}, \quad -\sqrt{\frac{a}{bk}}\,\frac{x}{y}, \quad +\sqrt{\frac{a}{bk^{2}}}\,\frac{x}{y}, \quad -\sqrt{\frac{a}{bk^{2}}}\,\frac{x}{y}.
\]
Hence substituting and observing that $ch(k-1)^{2} = (af + bgk^{2})$, it is easy to find
\[
\text{M} = c^{2}x^{2}k^{2}(k-1)^{4},
\]
that is, we have
\[
(bk^{2}\theta^{2} - a)(\theta y - x)^{2} + h(k-1)^{2}\theta^{2}w^{2} = 0.
\]

\newpage
\setcounter{page}{106}

If for greater convenience we write $\theta = \tfrac{1}{\phi}$, then this formula becomes
\[
(bk^{2} - a\phi^{2})(y - \phi x)^{2} + h(k-1)^{2}\phi^{2}w^{2} = 0,
\]
or, what is the same thing,
\[
(bk^{2}\phi^{2} - a)(y - \phi x)^{2} + h(k-1)^{2}\phi^{2}w^{2} = 0.
\]

Suppose that the terms in $U$ which contain $z^{2}$ are $= \Theta z^{2}$; then we have
\[
\Theta = \text{A}(1 - \phi_{1}^{2}y^{2})^{2}\,\Pi'\,[(bk^{2}\phi^{2} - a)(y - \phi x)^{2} + h(k-1)^{2}\phi^{2}w^{2}](k\phi - 1)^{2},
\]
where $\Pi'$ refers to the remaining three roots $\phi_{2}$, $\phi_{3}$, $\phi_{4}$; this may also be written
\[
\Theta = \frac{\text{A}\,\Pi\,[(bk^{2}\phi^{2} - a)(y - \phi x)^{2} + h(k-1)^{2}\phi^{2}w^{2}](k\phi - 1)^{2}}{(bk^{2}\phi_{1}^{2} - a)(y - \phi_{1}x)^{2} + h(k-1)^{2}\phi_{1}^{2}w^{2}}.
\]
Hence, observing that we have identically
\[
\Pi\,\{\phi w\sqrt{f} \pm iy\sqrt{g}\} = \mp\tfrac{c}{h}\{c\phi w\sqrt{f} \pm iy\sqrt{g}\}y^{2}\,\Pi'\,(\phi - \phi_{r}),
\]
and writing $\phi_{1} = \tfrac{1}{k}$, $h = \sqrt{-1}$ as usual, we find
%
\begin{multline}
\Pi\,(fx^{2}\phi^{2} + gy^{2})(k\phi - 1)^{2} = \tfrac{c^{2}}{h^{2}}\,\tfrac{1}{k^{2}}(bx^{2} - ay^{2} + hw^{2})^{2} \\
\times \Pi'\,(bk^{2}\phi^{2} - af + hv^{2});
\end{multline}
%
whence, writing
\[
\frac{\Pi\,(fx^{2}\phi^{2} + gy^{2})(k\phi - 1)^{2}}{\Pi\,[(bk^{2}\phi^{2} - a)(y - \phi x)^{2} + h(k-1)^{2}\phi^{2}w^{2}]} = \text{factor},
\]
we may calculate separately the terms A, B, and $2\Omega$ in
\[
\frac{\Theta}{x^{2}y^{2}} = \text{A}x^{4} + 2\Omega x^{2}y^{2} + \text{B}y^{4}.
\]

\newpage
\setcounter{page}{107}

The first of these is
%
\begin{multline}
\tfrac{c^{2}}{h^{2}}(k-1)^{2}(fx^{2} + k^{2}gy^{2})^{2}(bx^{2} - ay^{2} + hw^{2})^{2} \\
\div (fa^{2} + k^{2}gy^{2})^{2},
\end{multline}
%
if for shortness
%
\begin{multline}
(x, y, w)^{2} = (fx^{2} + gy^{2})[4\{(2-k)bx^{2} - ay^{2} + h(1-k)w^{2}\}^{2} \\
- 2(hx^{2} - ay^{2} + hw^{2})\{(6 - 6k + k^{2})bx^{2} - ay^{2} + h(1-k)^{2}w^{2}\}] \\
+ 4k^{2}(k-1)gy^{2}(bx^{2} - ay^{2} + hw^{2})\{(2-k)bx^{2} - ay^{2} + h(1-k)w^{2}\};
\end{multline}
%
the second is
\[
\tfrac{c^{2}}{h^{2}}(k-1)^{2}(fx^{2} + k^{2}gy^{2})^{2}(bx^{2} - ay^{2} + hw^{2})^{2} \div (fa^{2} + k^{2}gy^{2})^{2},
\]
if for shortness
%
\begin{multline}
(x, y, w)^{2} = \{(fx^{2} - k^{2}gy^{2})(cfx^{2} - cgy^{2} - fgw^{2}) - 2ckfgx^{2}y^{2}\} \\
\times [fgbk^{2}x^{2} + \{2ch(k-1)^{2} - af\}gy^{2} + fgh(k-1)^{2}w^{2}] \\
+ 2\{k^{2}bg - ch(k-1)^{2}\}fgx^{2}y^{2}[c(k+1)(fx^{2} - kgy^{2}) - kfgw^{2}];
\end{multline}
%
and hence
\[
\Theta = \frac{c^{2}(k-1)^{2}}{h^{2}(fx^{2} + k^{2}gy^{2})^{2}}[(x, y, w)_{1}^{2} + 2(bx^{2} - ay^{2} + hw^{2})(x, y, w)_{2}^{2}],
\]
which must be a rational and integral function of $(x, y, w)$.

In partial verification of this, observe that, because $U$ contains the terms $h^{2}cf(ch-af)x^{2}z^{2} + 2\Omega x^{2}y^{2}z^{2}w^{2}$,

\newpage
\setcounter{page}{108}

$\Theta$ should contain the terms
\[
2h^{2}cf(ch-af)x^{2} + 2\Omega x^{2}y^{2}w^{2},
\]
viz. in the term in $x^{2}$ should be $= 2h^{2}cf(ch-af)x^{2}$.

Now writing $y=0$, $w=0$, we have
\begin{align*}
\text{A} &= c^{2}f^{2}x^{4}, \\
(x, y, w)_{1}^{2} &= by\{4(2-k)^{2} - 2(6-6k+k^{2})\}x^{2}, \\
&= by(4 - 4k + 2k^{2})x^{2}, \\
(x, y, w)_{2}^{2} &= hcfg^{2}x^{2};
\end{align*}
and hence the required term of $\Theta$ is $x^{2}$ multiplied by
\[
\tfrac{c^{2}}{h^{2}}fh(4 - 4k + 2k^{2}) + 2h^{2}bcfg^{2}k^{2},
\]
viz. the coefficient is
\begin{align*}
&= 2b^{2}cf[ch(2-2k+k^{2}) + bgk^{2}] \\
&= 2b^{2}cf[ch + ch(1-k)^{2} + bgk^{2}],
\end{align*}
which in virtue of the relation $af + bgk^{2} + ch(1-k)^{2}$ becomes, as it should do,
\[
= 2b^{2}cf(ch-af).
\]

The actual division by $(fx^{2} + k^{2}gy^{2})^{2}$ would, however, be a very tedious process, and it is to be observed, that we only require to know the term $2\Omega x^{2}y^{2}w^{2}$ of $\Theta$.

We may therefore adopt a more simple course as follows: the terms of $\Theta$ which contain $w^{2}$ are $= (\text{A}x^{4} + 2\Omega x^{2}y^{2} + \text{B}y^{4})w^{2}$,

hence writing for a moment
\begin{align*}
(x, y, w)_{1}^{2} &= P + Qw^{2}, \\
(x, y, w)_{2}^{2} &= R + Sw^{2},
\end{align*}
and observing that we have
\begin{align*}
\text{A} &= c^{2}(fax^{2} + gy^{2} - 2c^{2}fg(fx^{2} - gy^{2})w^{2} + \&\text{c.}, \\
(bx^{2} - ay^{2} + hw^{2})^{2} &= (bx^{2} - ay^{2})^{2} + 2h(bx^{2} - ay^{2})w^{2} + \&\text{c.},
\end{align*}
we have
%
\begin{multline}
(fx^{2} + k^{2}gy^{2})(\text{A}x^{4} + 2\Omega x^{2}y^{2} + \text{B}y^{4}) = c^{2}h(fx^{2} + gy^{2})Q - 2c^{2}fgh(fx^{2} - gy^{2})P \\
+ (bx^{2} - ay^{2})^{2}S + 2h(bx^{2} - ay^{2})R.
\end{multline}
%
But in this identical equation we may write $x^{2}=a$, $y^{2}=b$, which gives
\[
(af + k^{2}bg)(\text{A}x^{4} + 2\Omega ab + \text{B}y^{4}) = c^{2}h(af + bg)Q - 2c^{2}fgh(af - bg)P;
\]
and from the equation
\[
(x, y, w)_{1}^{2} = P + Qw^{2},
\]
we have
%
\begin{multline}
P + Qw^{2} = (af + bg)^{2}\biggl[\frac{(1-k)^{2}ab + (1-k)hw^{2}}{\cdots}\biggr] \\
+ 4k^{2}(k-1)bghw^{2}(1-k)ab.
\end{multline}

\newpage
\setcounter{page}{109}

that is,
\[
Q = (k-1)^{2}abh[ch(-6k^{2} + 4k + 2) - 4k^{2}bg],
\]
whence
%
\begin{multline}
(k-1)^{2}(\text{A}x^{4} + 2\Omega ab + \text{B}y^{4}) = ab(af + bg)^{2}[ch(-6k^{2} + 4k + 2) - 4k^{2}bg] \\
+ 8fg(af - bg)(k-1)^{2}(ab)^{2}.
\end{multline}
%
But we have
\[
\text{A}x^{4} + \text{B}y^{4} = -2ab(af - bg)(-afbg + c^{2}h^{2} + 2chaf + 2chbg),
\]
and thence
%
\begin{multline}
2(k-1)^{2}\Omega = (af + bg)^{2}[ch(-6k^{2} + 4k + 2) - 4k^{2}bg] \\
+ (k-1)^{2}(af - bg)\{-2afbg + 2c^{2}h^{2} + 4chaf + 4chbg\},
\end{multline}
%
or
%
\begin{multline}
(k-1)^{2}\Omega = (af + bg)^{2}[ch(-3k^{2} + 2k + 1) - 2k^{2}bg] \\
+ (k-1)^{2}(af - bg)\{3afbg + c^{2}h^{2} + 2chaf + 2chbg\}.
\end{multline}
%
Writing $-3k^{2} + 2k + 1 = -3(k-1)^{2} - 4(k-1)$, this is
\[
\Omega = (af + bg)^{2}[ch(-3 - \tfrac{4}{k-1}) - \tfrac{2k^{2}}{(k-1)^{2}}bg] + (af - bg)\{3afbg + c^{2}h^{2} + 2chaf + 2chbg\},
\]
or since
\[
1 : -k : k-1 = \lambda : \mu : \nu;
\]
and writing for shortness
\[
(af,\; bg,\; ch) = (\alpha,\; \beta,\; \gamma),
\]
this is
\[
\Omega = (\alpha + \beta)^{2}\gamma(-3 - \tfrac{4\lambda}{\nu}) + (\alpha - \beta)\{3\alpha\beta + \gamma^{2} + 2\gamma\alpha + 2\gamma\beta\},
\]
which is the value of $\Omega$: viz. the conclusion arrived at is that, eliminating $\theta$ from the equations
\begin{align*}
(h\theta^{2} + g)(k\theta y - x)^{2} + h(k-1)^{2}\theta^{2}z^{2} &= 0, \\
(bk^{2}\theta^{2} - a)(\theta y - x)^{2} + h(k-1)^{2}\theta^{2}w^{2} &= 0,
\end{align*}
where $k$ denotes a determinate function of $af$, $bg$, $ch$, viz. writing $af$, $bg$, $ch = \alpha$, $\beta$, $\gamma$ and
\[
1 : -k : k-1 = \lambda : \mu : \nu,
\]
we have
\[
\lambda + \mu + \nu = 0, \qquad \alpha\lambda^{2} + \beta\mu^{2} + \gamma\nu^{2} = 0,
\]
equations which serve to determine $k$: the result of the elimination is the octic equation
\[
b^{2}c^{2}f^{2}x^{4} + \ldots + 2\Omega x^{2}y^{2}z^{2}w^{2} = 0,
\]
where $\Omega$ has the last-mentioned value.

\medskip\noindent\textsc{C.~X.}\hfill 12

\newpage
\setcounter{page}{110}

The value of $\Omega$ is unsymmetrical in its form, and there are apparently six values; viz. writing
\begin{align*}
\text{A} &= (\beta + \gamma)^{2}\alpha(-3 - \tfrac{4\lambda}{\nu}) + (\alpha - \beta)(-3 + \alpha\beta + \gamma^{2}), \\
\text{B} &= (\gamma + \alpha)^{2}\beta(-3 - \tfrac{4\mu}{\lambda}) + (\beta - \gamma)(-3 + \beta\gamma + \alpha^{2}), \\
\text{C} &= (\alpha + \beta)^{2}\gamma(-3 - \tfrac{4\nu}{\mu}) + (\gamma - \alpha)(-3 + \gamma\alpha + \beta^{2}), \\
\text{A}_{1} &= (\beta + \gamma)^{2}\alpha(-3 - \tfrac{4\lambda}{\nu}) - (\alpha - \beta)(-3 + \alpha\beta + \gamma^{2}), \\
\text{B}_{1} &= (\gamma + \alpha)^{2}\beta(-3 - \tfrac{4\mu}{\lambda}) - (\beta - \gamma)(-3 + \beta\gamma + \alpha^{2}), \\
\text{C}_{1} &= (\alpha + \beta)^{2}\gamma(-3 - \tfrac{4\nu}{\mu}) - (\gamma - \alpha)(-3 + \gamma\alpha + \beta^{2}),
\end{align*}
where for shortness $S = 2(\beta\gamma + \gamma\alpha + \alpha\beta)$, the six values would be A, B, C, A$_{1}$, B$_{1}$, C$_{1}$.

But we have really
\[
\text{A} = \text{B} = \text{C} = -\text{A}_{1} = -\text{B}_{1} = -\text{C}_{1},
\]
so that $\Omega$ has really only two values, equal and of opposite signs, or, what is the same thing, $\Omega^{2}$ has a unique value.

In fact, writing for shortness
\[
\lambda + \mu + \nu = P, \qquad \alpha\lambda^{2} + \beta\mu^{2} + \gamma\nu^{2} = Z,
\]
we find at once the identity
\[
\lambda^{2}(\text{A} + \text{A}_{1}) = (\beta + \gamma)^{2}\{-2Z - 4\lambda\alpha P\},
\]
so that $\text{A} = -\text{A}_{1}$ in virtue of $P=0$, $Z=0$.

And similarly B~$= -\text{B}_{1}$, C~$= -\text{C}_{1}$.

But the demonstration of the equation $\text{A} = \text{B}$ is more complicated. We have
%
\begin{multline}
\text{A} - \text{B} = -3\alpha(\beta + \gamma)^{2} + 3\beta(\gamma + \alpha)^{2} \\
- 4\alpha(\beta + \gamma)^{2}\tfrac{\mu}{\nu} - 2\gamma(\beta + \gamma)^{2}\tfrac{\lambda^{2}}{\mu^{2}} \\
+ 4\beta(\gamma + \alpha)^{2}\tfrac{\nu}{\lambda} + 2\alpha(\gamma + \alpha)^{2}\tfrac{\mu^{2}}{\nu^{2}} \\
+ (\beta - \gamma)(S + \beta\gamma + \alpha^{2}) - (\gamma - \alpha)(S + \gamma\alpha + \beta^{2}),
\end{multline}
%
that is,
%
\begin{multline}
\lambda^{2}\nu^{2}(\text{A} - \text{B}) = \{-3\alpha(\beta + \gamma)^{2} + 3\beta(\gamma + \alpha)^{2} + (\beta - \gamma)(S + \beta\gamma + \alpha^{2}) - (\gamma - \alpha)(S + \gamma\alpha + \beta^{2})\}\lambda^{2}\nu^{2} \\
- 4\alpha(\beta + \gamma)^{2}\lambda\mu\nu^{2} \\
- 2\gamma(\beta + \gamma)^{2}\nu^{2}\lambda^{2} \\
+ 4\beta(\gamma + \alpha)^{2}\lambda^{2}\mu\nu \\
+ 2\alpha(\gamma + \alpha)^{2}\lambda^{2}\mu^{2}.
\end{multline}

\newpage
\setcounter{page}{111}

Or, denoting for a moment the coefficient of $\lambda^{2}\nu^{2}$ by K, and writing also $\nu^{2} = \lambda - \alpha\lambda^{2} - \beta\mu^{2}$, $\nu = P - \lambda - \mu$, this is
%
\begin{multline}
= K\lambda^{2}\nu^{2} - 4\alpha(\beta + \gamma)^{2}\lambda\mu^{2} \\
- 2(\beta + \gamma)^{2}\mu^{2}(\lambda - \alpha\lambda^{2} - \beta\mu^{2}) \\
+ 4\beta(\gamma + \alpha)^{2}\lambda\nu(P - \lambda - \mu) \\
= -2(\beta + \gamma)^{2}\mu^{2}\lambda + 4\beta(\gamma + \alpha)^{2}\lambda^{2}P \\
+ 2\alpha(\gamma + \alpha)^{2}\lambda^{2} \\
- 4\beta(\gamma + \alpha)^{2}\lambda^{2}\mu \\
+ \{-4\alpha(\beta + \gamma)^{2} + \cdots\}\lambda\mu^{2} \\
+ 2\beta(\beta + \gamma)^{2}\mu^{4},
\end{multline}
%
and here the coefficient of $\lambda^{2}\mu^{2}$ is found to be
\[
= 2\{\alpha\beta S(\alpha + \beta) + \gamma(\alpha - \beta)^{2} - 3\gamma^{2}(\alpha + \beta)\}.
\]
Hence, the terms without $\lambda$ or $P$ are $= 2\text{V}$, where
%
\begin{multline}
\text{V} = \alpha(\gamma + \alpha)^{2}\lambda^{4} - 2\beta(\gamma + \alpha)^{2}\lambda^{3}\mu \\
+ \{\alpha\beta S(\alpha + \beta) + \gamma(\alpha - \beta)^{2} - 3\gamma^{2}(\alpha + \beta)\}\lambda^{2}\mu^{2} \\
- 2\alpha(\beta + \gamma)^{2}\lambda\mu^{3} + \beta(\beta + \gamma)^{2}\mu^{4},
\end{multline}
%
and this is identically
%
\begin{multline}
(\alpha + \gamma)\text{V} = [\alpha(\gamma + \alpha)\lambda^{2} - (\beta\gamma + \gamma\alpha + \alpha\beta)\lambda\mu + \beta(\beta + \gamma)\mu^{2}]^{2} \\
- \{\alpha\lambda^{2} + \beta\mu^{2} + \gamma(P - \lambda - \mu)^{2}\}\times\{
\alpha(\gamma + \alpha)^{2}\lambda^{2} - 2(\beta\gamma + \gamma\alpha + \alpha\beta)\lambda\mu + \beta(\beta + \gamma)\mu^{2}\},
\end{multline}
%
where observing that
\[
\alpha\lambda^{2} + \beta\mu^{2} + \gamma\nu^{2} = Z,
\]
we have the first factor
\[
(\alpha + \gamma)\lambda^{2} + (\beta + \gamma)\mu^{2} + 2\gamma\nu\lambda = -\lambda - \gamma\nu^{2} + 2\gamma P(\lambda + \mu^{2}),
\]
and consequently
%
\begin{multline}
\lambda^{2}\nu^{2}(\text{A} - \text{B}) = -2(\beta + \gamma)^{2}\mu^{2}Z + 4\beta(\gamma + \alpha)^{2}\lambda^{2}P \\
+ 2\{Z - \gamma\nu^{2} + 2\gamma P(\lambda + \mu)\}[\alpha(\gamma + \alpha)\lambda^{2} - 2(\beta\gamma + \gamma\alpha + \alpha\beta)\lambda\mu + \beta(\beta + \gamma)\mu^{2}]:
\end{multline}
%
viz. in virtue of $P=0$, $Z=0$, we have $\text{A} = \text{B}$.

And thus
\[
\text{A} = \text{B} = \text{C} = -\text{A}_{1} = -\text{B}_{1} = -\text{C}_{1}:
\]
so that the only values of $\Omega$ are, say, A and~$-A$.

\medskip\noindent 12---2

\newpage
\setcounter{page}{112}

Reverting to the original equations
\begin{align*}
(h\theta^{2} + g)(k\theta y - x)^{2} + h(k-1)^{2}\theta^{2}z^{2} &= 0, \\
(bk^{2}\theta^{2} - a)(\theta y - x)^{2} + h(k-1)^{2}\theta^{2}w^{2} &= 0,
\end{align*}
say these are
\[
(a, b, c, d, e\,\!\!\between\,\theta, y)^{2} = 0, \qquad (a', b', c', d', e'\,\!\!\between\,\theta, y)^{2} = 0,
\]
then the coefficients in the two equations have the values
\begin{align*}
a &= f\theta^{2}y^{2}, & a' &= bk^{2}\theta^{2}, \\
b &= -2kf\theta xy, & b' &= -2bk^{2}\theta^{2}xy, \\
c &= fx^{2} + gk^{2}y^{2} + h(k-1)^{2}z^{2}, & c' &= hk^{2}x^{2} - ay^{2} + h(k-1)^{2}w^{2}, \\
d &= -2gk\theta xy, & d' &= 2a\theta xy, \\
e &= g\theta^{2}y^{2}, & e' &= -a\theta^{2}y^{2},
\end{align*}
where observe that only $c$ contains $z^{2}$, and only $c'$ contains $w^{2}$.

The result of the elimination is
\[
\begin{vmatrix}
a & b & c & d & e \\
  & a & b & c & d \\
  &   & a & b & c \\
a' & b' & c' & d' & e' \\
   & a' & b' & c' & d' \\
   &    & a' & b' & c'
\end{vmatrix}
= 0;
\]
viz. here the only terms which contain $z^{2}$ and $w^{2}$ are
\[
c^{3}w^{2} \cdot c'^{3}z^{2},
\]
and hence the terms in $z^{4}$ and $w^{4}$ are
\begin{align*}
&h^{2}(k-1)^{4}z^{4} \cdot a^{2}b'k^{2}x^{2}y^{2} + h^{2}(k-1)^{4}w^{4} \cdot f^{2}g^{2}k^{2}x^{2}y^{2}, \\
&= h^{2}k^{2}(k-1)^{4}x^{2}y^{2}(a^{2}b' \cdot z^{4} + f^{2}g^{2} \cdot w^{4}),
\end{align*}
viz. these are
\[
\text{or assuming that the determinant contains as a factor the function}
\]
\[
b^{2}c^{2}f^{2}w^{2}x^{2} + \ldots + 2\Omega x^{2}y^{2}z^{2}w^{2},
\]
with a properly determined value of $\Omega$, we see that the other factor is $= h^{2}k^{2}(k-1)^{4}x^{2}y^{2}$, which agrees with a preceding result.

\newpage

% #############################################################################
% PAPER 655: A MEMOIR ON DIFFERENTIAL EQUATIONS
% From Quarterly Journal of Pure and Applied Mathematics, vol. XIV (1877), pp. 292--339
% Original article pages 93--130
% #############################################################################
\section*{655. A Memoir on Differential Equations.}
\addcontentsline{toc}{section}{655. A Memoir on Differential Equations}

\begin{center}
[From the \textit{Quarterly Journal of Pure and Applied Mathematics,} vol.~\textsc{xiv}. (1877), pp.~292--339.]
\end{center}

\setcounter{page}{113}

We have to do with a set of variables, which is either unipartite $(x, y, z, \ldots)$, or else bipartite $(x, y, z, \ldots;\; p, q, r, \ldots)$, the variables in the latter case corresponding in pairs $x$~and~$p$, $y$~and~$q$, \&c.

A letter not otherwise explained denotes a function of the variables. Any such letter may be put $= \text{const.}$, viz.~we thereby establish a relation between the variables and when this is so, we use the same letter to denote the constant value of the function.

Thus the set being $(x, y, z;\; p, q, r)$, $H$ may denote a given function $pqr - xyz$; and then, if $H = \text{const.}$, we have $pqr - xyz = H$ (a constant). This notation, when once clearly understood, is, I think, a very convenient one.

The present memoir relates chiefly to the following subjects:

\textbf{A.} Unipartite set $(x, y, z, \ldots)$. The differential system
\[
\frac{dx}{X} = \frac{dy}{Y} = \frac{dz}{Z} = \cdots
\]
and connected therewith the linear partial differential equation
\[
X\frac{df}{dx} + Y\frac{df}{dy} + Z\frac{df}{dz} + \cdots = 0:
\]
also the lineo-differential $X\,dx + Y\,dy + Z\,dz + \cdots$.

\textbf{B.} Bipartite set $(x, y, z, \ldots;\; p, q, r, \ldots)$. The Hamiltonian system
\[
\frac{dx}{dp} = \frac{dy}{dq} = \frac{dz}{dr} = \cdots = \frac{dp}{dH} = \frac{dq}{dH} = \frac{dr}{dH},
\]

\newpage
\setcounter{page}{114}

and connected therewith the linear partial differential equation
\[
\frac{dH}{dp}\frac{d\Theta}{dx} - \frac{dH}{dx}\frac{d\Theta}{dp} + \frac{dH}{dq}\frac{d\Theta}{dy} - \frac{dH}{dy}\frac{d\Theta}{dq} + \cdots = 0,
\]
otherwise written
\[
\frac{d(H, \Theta)}{d(p, x)} + \frac{d(H, \Theta)}{d(q, y)} + \frac{d(H, \Theta)}{d(r, z)} + \cdots = 0,
\]
where $H$ denotes a given function of the variables: also the Hamiltonian system as augmented by an equality $= dt$, and as augmented by this and another equality
\begin{align*}
dH &= \frac{dH}{dp}\,dp + \frac{dH}{dq}\,dq + \frac{dH}{dr}\,dr + \cdots \\
&= -\biggl(\frac{dH}{dx}\,dx + \frac{dH}{dy}\,dy + \frac{dH}{dz}\,dz + \cdots\biggr).
\end{align*}

\textbf{C.} Bipartite set $(x, y, z, \ldots;\; p, q, r, \ldots)$. The partial differential equation $H = \text{const.}$, where, as before, $H$ is a given function of the variables, but $p$, $q$, $r$, \ldots\ are now the differential coefficients in regard to $x$, $y$, $z$, \ldots\ respectively of a function $V$ of these variables, or, what is the same thing, there exists a function
\[
V = \int(p\,dx + q\,dy + r\,dz + \cdots),
\]
of the variables $x$, $y$, $z$, \ldots

In what precedes, I have written $(x, y, z, \ldots)$ to denote a set of any number~$n$ of variables, and $(x, y, z, \ldots;\; p, q, r, \ldots)$ to denote a set of any even number~$2n$ of variables, and the investigations are for the most part applicable to these general cases. But for greater clearness and facility of expression, I usually consider the case of a set $(x, y, z, w)$, or $(x, y, z;\; p, q, r)$, \&c., as the case may be, consisting of a definite number of variables.

The greater part of the theory is not new, but I think that I have presented it in a more compact and intelligible form than has hitherto been done, and I have added some new results.

\subsection*{Introductory Remarks. Art.~Nos.~1 to~3.}

\subsubsection*{1.}

As already noticed, a letter not otherwise explained is considered as denoting a function of the variables of the set; but when necessary we indicate the variables by a notation such as $z = z(x, y)$; $z$ is here a function (known or unknown as the case may be) of the variables $x$, $y$, the $z$ on the right-hand side being in fact a functional symbol.

And thus also $z = z(x, y) = \text{const.}$ denotes that the function $z(x, y)$ of the variables $x$, $y$ has a constant value, which constant value is $=z$, viz.~we thus indicate a relation between the variables $x$,~$y$.

\newpage
\setcounter{page}{115}

\subsubsection*{2.}

The variables $x$, $y$, \&c., may have infinitesimal increments $dx$, $dy$, \&c.; and the equations of connexion between the variables then give rise to linear relations between these increments, the coefficients therein being differential coefficients and, as such, represented in the usual notation; thus if $z = z(x, y)$, we have
\[
dz = \frac{dz}{dx}\,dx + \frac{dz}{dy}\,dy,
\]
where $\dfrac{dz}{dx}$, $\dfrac{dz}{dy}$ are the so-called partial differential coefficients of $z$ in regard to $x$, $y$ respectively.

If we have $y = y(x)$, then also $dy = \dfrac{dy}{dx}\,dx$, and the foregoing equation becomes
\[
\frac{d(z)}{dx} = \frac{dz}{dx} + \frac{dz}{dy}\,\frac{dy}{dx};
\]
but considering the two equations $z = z(x, y)$ and $y = y(x)$ as determining $z$ as a function of $x$, say $z = z(x)$, we have $dz = \dfrac{d(z)}{dx}\,dx$; whence comparing the two formulae
\[
\frac{d(z)}{dx} = \frac{dz}{dx} + \frac{dz}{dy}\,\frac{dy}{dx}.
\]
The distinction is best made, not by any difference of notation $\dfrac{dz}{dx}$, $\dfrac{d(z)}{dx}$, but by appending in any case of doubt the equations or equation used in the differentiation. Thus we have $\dfrac{dz}{dx}$, where $z = z(x, y)$: or, as the case may be, $\dfrac{d(z)}{dx}$ where $z = z(x, y)$ and $y = y(x)$.

\subsubsection*{3.}

A relation between increments is always really a relation between differential coefficients: but we use the increments for symmetry and conciseness, as in the case of a differential system $\dfrac{dx}{X} = \dfrac{dy}{Y} = \dfrac{dz}{Z}$, or in a question relating to the lineo-differential $X\,dx + Y\,dy + Z\,dz$, for instance in the question whether this can be put $= du$.

\subsection*{Notations. Art.~Nos.~4 to~6.}

\subsubsection*{4. Functional determinants.}

If $a$, $b$, $c$, \ldots\ are functions of the variables $x$, $y$, $z$, $w$, \ldots, then the determinants
\[
\begin{vmatrix}
\dfrac{da}{dx} & \dfrac{da}{dy} \\[0.8em]
\dfrac{db}{dx} & \dfrac{db}{dy}
\end{vmatrix},
\quad
\begin{vmatrix}
\dfrac{da}{dx} & \dfrac{da}{dy} & \dfrac{da}{dz} \\[0.5em]
\dfrac{db}{dx} & \dfrac{db}{dy} & \dfrac{db}{dz} \\[0.5em]
\dfrac{dc}{dx} & \dfrac{dc}{dy} & \dfrac{dc}{dz}
\end{vmatrix},
\quad \&\text{c.},
\]

\newpage
\setcounter{page}{116}

are for shortness represented by
\[
\frac{d(a, b)}{d(x, y)}, \qquad \frac{d(a, b, c)}{d(x, y, z)}, \qquad \&\text{c.},
\]
the notation being especially used in the first-mentioned case where the symbol is $= \left(\dfrac{a}{x}\right)$.

It is sometimes convenient to extend this notation, and for instance use $\dfrac{d(a, b)}{d(x, y, z)}$ to denote the series of determinants
\[
\begin{vmatrix}
\dfrac{da}{dx} & \dfrac{da}{dy} & \dfrac{da}{dz} \\[0.5em]
\dfrac{db}{dx} & \dfrac{db}{dy} & \dfrac{db}{dz}
\end{vmatrix},
\]
which can be formed by selecting in every way two columns to form thereout a determinant; the equation $\dfrac{d(a, b)}{d(x, y, z)} = 0$ will then denote that each of these determinants is $= 0$.

The analogous notation $\dfrac{d(a, b, c)}{d(x, y)}$ would denote non-existent determinants, viz.~there are here not columns enough to form with them a determinant: and the notation is not required.

\subsubsection*{5.}

In the case of a bipartite set $(x, y, z, \ldots;\; p, q, r, \ldots)$, if $a$, $b$ are any functions of these variables, we consider the derivative
\[
(a, b) = \frac{d(a, b)}{d(p, x)} + \frac{d(a, b)}{d(q, y)} + \frac{d(a, b)}{d(r, z)} + \cdots,
\]
viz.~$(a, b)$ is used to denote the sum of the functional determinants on the right hand.

\subsubsection*{6.}

Taking again $(x, y, z, w, \ldots)$ as the variables, then in the theory of the lineo-differential $X\,dx + Y\,dy + Z\,dz + W\,dw + \cdots$, we use certain derivative functions analogous to Pfaffians.

They may be thus defined; viz.~considering the numbers $1$, $2$, $3$, $4$, \ldots\ as corresponding to the variables $x$, $y$, $z$, $w$, \ldots\ respectively, we have
\begin{align*}
1 &= X, \quad 2 = Y, \quad 3 = Z, \quad 4 = W, \quad \&\text{c.}, \\
12 &= \frac{dX}{dy} - \frac{dY}{dx}, \quad
13 = \frac{dX}{dz} - \frac{dZ}{dx}, \quad
23 = \frac{dY}{dz} - \frac{dZ}{dy}, \quad \&\text{c.}, \\
123 &= 1.23 + 2.31 + 3.12 \\
&= X\biggl(\frac{dY}{dz} - \frac{dZ}{dy}\biggr) + Y\biggl(\frac{dZ}{dx} - \frac{dX}{dz}\biggr) + Z\biggl(\frac{dX}{dy} - \frac{dY}{dx}\biggr), \\
1234 &= 12.34 + 13.42 + 14.23 \\
&= \biggl(\frac{dX}{dy} - \frac{dY}{dx}\biggr)\biggl(\frac{dZ}{dw} - \frac{dW}{dz}\biggr) + \biggl(\frac{dX}{dz} - \frac{dZ}{dx}\biggr)\biggl(\frac{dW}{dy} - \frac{dY}{dw}\biggr) \\
&\qquad + \biggl(\frac{dX}{dw} - \frac{dW}{dx}\biggr)\biggl(\frac{dY}{dz} - \frac{dZ}{dy}\biggr).
\end{align*}

\newpage
\setcounter{page}{117}

and, adding for greater distinctness the next following cases,
\begin{align*}
12345 &= 1.2345 + 2.3451 + 3.4512 + 4.5123 + 5.1234, \\
123456 &= 12.3456 + 13.4561 + 14.5612 + 15.6123 + 16.2345,
\end{align*}
where of course $2345$, \&c., have the significations mentioned above.

\subsection*{Dependency of Functions. Art.~Nos.~7 and~8.}

\subsubsection*{7.}

Two or more functions of the same variables may be independent, or else dependent or connected; viz.~in the latter case any one of the functions is a function of the others. If $a = a(x)$, $b = b(x)$, the functions $a$, $b$ are dependent, but if $a = a(x, y)$, $b = b(x, y)$, then the condition of dependency is
\[
\frac{d(a, b)}{d(x, y)} = 0,
\]
and, similarly, if $a = a(x, y, z)$, $b = b(x, y, z)$, then the conditions of dependency are
\[
\frac{d(a, b)}{d(x, y, z)} = 0,
\]
viz.~if the equations thus represented are all of them satisfied, the functions are dependent, but if not, then they are independent.

Observe that, when $a = a(x, y, z)$, $b = b(x, y, z)$ as above, if we choose to attend only to the variables $x$, $y$, treating $z$ as a mere constant, there is then a single condition of dependency $\dfrac{d(a, b)}{d(x, y)} = 0$, and so if we attend only to the variable $x$, treating $y$, $z$ as mere constants, then $a$ and $b$ are dependent.

Thus when $a = x$, $b = x^{2} + y$, the functions $a$, $b$ are independent if we attend to both the variables $x$, $y$; dependent if $y$ be regarded as a constant.

\subsubsection*{8.}

Further when $a = a(x, y)$, $b = b(x, y)$, $c = c(x, y)$, the functions $a$, $b$, $c$ are dependent; but when $a = a(x, y, z)$, $b = b(x, y, z)$, $c = c(x, y, z)$, the condition of dependency is
\[
\frac{d(a, b, c)}{d(x, y, z)} = 0,
\]
and so when $a = a(x, y, z, w)$, $b = b(x, y, z, w)$, $c = c(x, y, z, w)$, the conditions of dependency are
\[
\frac{d(a, b, c)}{d(x, y, z, w)} = 0,
\]
viz.~if all the equations thus represented are satisfied, the functions are dependent; but if not, then they are independent.

And so in other cases.

\medskip\noindent\textsc{C.~X.}\hfill 13

\newpage
\setcounter{page}{118}

\subsection*{The General Differential System. Art.~Nos.~9 to~22.}

\subsubsection*{9.}

Taking the set of variables to be $(x, y, z, w)$, the system is
\[
\frac{dx}{X} = \frac{dy}{Y} = \frac{dz}{Z} = \frac{dw}{W},
\]
and we associate with this the linear partial differential equation
\[
X\frac{df}{dx} + Y\frac{df}{dy} + Z\frac{df}{dz} + W\frac{df}{dw} = 0.
\]

\subsubsection*{10.}

It is tolerably evident that the differential equations establish between $x$, $y$, $z$, $w$ a threefold relation depending upon three arbitrary constants; in fact, regarding $(x, y, z, w)$ as the coordinates of a point in four-dimensional space, and starting from any given point, the differential equations determine the ratios of the increments $dx$, $dy$, $dz$, $dw$, that is, the direction of passage to a consecutive point and then again taking for $x$, $y$, $z$, $w$ the coordinates of this point, the same equations give the direction of passage to the next consecutive point, and so on.

The locus of the point is therefore a curve, or we have between the coordinates a threefold relation, and (the initial point being arbitrary) we have a curve of the system through each point of the four-dimensional space, viz.~the relation must involve three arbitrary constants.

But this being so, the constants will be expressible as functions of the coordinates, viz.~the threefold relation involving the three constants will be expressible in the form $a = \text{const.}$, $b = \text{const.}$, $c = \text{const.}$, where $a$, $b$, $c$ denote respective functions of $(x, y, z, w)$; and the differential equations thus establish between the variables the threefold relation $a = \text{const.}$, $b = \text{const.}$, $c = \text{const.}$.

\newpage
\setcounter{page}{119}

\subsubsection*{11.}

Observe that, in speaking of an integral $a$, we consider it as a function of the variables such that the differential equation gives $da = 0$; viz.~$a$ being a function of $(x, y, z, w)$, if its differential is
\[
da = A\,dx + B\,dy + C\,dz + D\,dw,
\]
the condition is
\[
AX + BY + CZ + DW = 0.
\]
This equation must be satisfied identically, and we have therefore a single partial differential equation
\[
X\frac{da}{dx} + Y\frac{da}{dy} + Z\frac{da}{dz} + W\frac{da}{dw} = 0,
\]
which must be satisfied by the function~$a$; the general integral of this partial differential equation contains an arbitrary function, and we may say that $a$ is any function whatever of $(x, y, z, w)$ which satisfies this partial differential equation.

Similarly $b$ is any other function which satisfies the same partial differential equation, and so for~$c$.

The number of integrals is thus unlimited, but if $a$ be any particular integral, then every other integral is a function of $(a, b, c)$.

\subsubsection*{12.}

In particular, if $X$, $Y$, $Z$, $W$ are linear functions of the variables, say
\begin{align*}
X &= a_{1}x + b_{1}y + c_{1}z + d_{1}w, \\
Y &= a_{2}x + b_{2}y + c_{2}z + d_{2}w, \\
Z &= a_{3}x + b_{3}y + c_{3}z + d_{3}w, \\
W &= a_{4}x + b_{4}y + c_{4}z + d_{4}w,
\end{align*}
then a particular integral is had by assuming $a = lx + my + nz + pw$, where $l$, $m$, $n$, $p$ are constants; and the condition $da = 0$ gives at once
\begin{align*}
l(a_{1}x + b_{1}y + c_{1}z + d_{1}w) + m(a_{2}x + b_{2}y + c_{2}z + d_{2}w) \\
+ n(a_{3}x + b_{3}y + c_{3}z + d_{3}w) + p(a_{4}x + b_{4}y + c_{4}z + d_{4}w) &= 0,
\end{align*}
that is,
\begin{align*}
l a_{1} + m a_{2} + n a_{3} + p a_{4} &= 0, \\
l b_{1} + m b_{2} + n b_{3} + p b_{4} &= 0, \\
l c_{1} + m c_{2} + n c_{3} + p c_{4} &= 0, \\
l d_{1} + m d_{2} + n d_{3} + p d_{4} &= 0,
\end{align*}
equations which determine the ratios $(l : m : n : p)$.

\newpage
\setcounter{page}{120}

\subsubsection*{13.}

The general integral $a = \text{const.}$ is in fact any arbitrary function of the variables $(x, y, z, w)$ which satisfies the partial differential equation
\[
X\frac{da}{dx} + Y\frac{da}{dy} + Z\frac{da}{dz} + W\frac{da}{dw} = 0.
\]
Taking $a$ to be any particular integral, then every other integral is a function of $a$; for if $b$ be any other integral, then $b$ is a function of $(x, y, z, w)$ which satisfies the same partial differential equation, and we have therefore $b = \phi(a)$.

The general integral may be written in the form
\[
\Phi(a, b, c) = 0,
\]
where $\Phi$ is an arbitrary functional symbol; and this is equivalent to the system of three equations $a = \text{const.}$, $b = \text{const.}$, $c = \text{const.}$, involving as it should do three arbitrary constants.

\subsubsection*{14.}

If in the differential system we have two integrals $a$, $b$, these may be regarded as given functions of $(x, y, z, w)$; and we may then express the ratios $dx : dy : dz : dw$ by means of the Jacobians. In fact, regarding $a$, $b$ as given functions, we have
\[
\frac{dx}{\,d(a, b)/d(y, z, w)\,} = \frac{dy}{\,d(a, b)/d(z, w, x)\,} = \frac{dz}{\,d(a, b)/d(w, x, y)\,} = \frac{dw}{\,d(a, b)/d(x, y, z)\,},
\]
where the notation is used to denote the series of functional determinants
\[
\frac{d(a, b)}{d(y, z)}, \quad \frac{d(a, b)}{d(z, w)}, \quad \&\text{c.},
\]
and similarly for the other denominators.

\subsubsection*{15.}

The foregoing ratios of $dx : dy : dz : dw$ may also be written in the form
\[
\frac{dx}{\,\partial(a, b)/\partial(y, z)\,} = \frac{dy}{\,\partial(a, b)/\partial(z, x)\,} = \frac{dz}{\,\partial(a, b)/\partial(x, y)\,} = \frac{dw}{\,\partial(a, b)/\partial(x, y, z)\,},
\]
where the notation is used to denote the series of partial derivatives
\[
\frac{\partial a}{\partial y}\frac{\partial b}{\partial z} - \frac{\partial a}{\partial z}\frac{\partial b}{\partial y}, \quad \&\text{c.},
\]
and similarly for the other denominators.

\newpage
\setcounter{page}{121}

\subsubsection*{16.}

There is no general process for the integration of the differential system
\[
\frac{dx}{X} = \frac{dy}{Y} = \frac{dz}{Z} = \frac{dw}{W},
\]
but in particular cases the integration may be effected by various methods; for instance, if we can find a factor $\lambda$ such that
\[
\lambda(X\,dx + Y\,dy + Z\,dz + W\,dw) = du,
\]
where $u$ is a function of $(x, y, z, w)$, then $u = \text{const.}$ is an integral.

\subsubsection*{17.}

In particular, if the expression $X\,dx + Y\,dy + Z\,dz + W\,dw$ is an exact differential, say $= du$, then $u = \text{const.}$ is an integral; and if the expression is a perfect differential, say $= du + v\,dw$, then $u + vw = \text{const.}$ is an integral.

\subsubsection*{18.}

The condition that $X\,dx + Y\,dy + Z\,dz + W\,dw$ should be an exact differential is that the Pfaffian derivatives should all vanish; viz.~that we should have
\[
12 = 0, \quad 13 = 0, \quad 23 = 0, \quad 14 = 0, \quad 24 = 0, \quad 34 = 0,
\]
or, what is the same thing,
\[
\frac{dX}{dy} = \frac{dY}{dx}, \quad \frac{dX}{dz} = \frac{dZ}{dx}, \quad \frac{dY}{dz} = \frac{dZ}{dy}, \quad \&\text{c.}
\]

\subsubsection*{19.}

The result would have been obtained by assuming $a$, $b$, $c$ as known functions of $(x, y, z, w)$; say, for example, $a = x$, $b = y$, and seeking to express $c$ as a function of $(x, y, z, w)$, or say of $(x, y, z)$, for the determination of the ratios $dx : dy : dz : dw$.

\subsubsection*{20.}

The differential system may be treated by the method of multipliers; if $\lambda$, $\mu$, $\nu$, $\rho$ be factors such that
\[
\lambda\,dx + \mu\,dy + \nu\,dz + \rho\,dw = 0,
\]
then we have
\[
\lambda X + \mu Y + \nu Z + \rho W = 0,
\]
and the condition that this should be an exact differential gives rise to the partial differential equations for the determination of $\lambda$, $\mu$, $\nu$, $\rho$.

\newpage
\setcounter{page}{122}

\subsubsection*{21.}

The method of multipliers may be applied as follows: assuming that we have two integrals $a$, $b$, then the differential system may be written
\[
\frac{dx}{\partial(a, b)/\partial(y, z, w)} = \frac{dy}{\partial(a, b)/\partial(z, w, x)} = \cdots,
\]
where the denominators are the Jacobian determinants; and the condition for a third integral is that the expression
\[
\frac{\partial(a, b)}{\partial(y, z, w)}\,dx + \frac{\partial(a, b)}{\partial(z, w, x)}\,dy + \cdots
\]
should be an exact differential.

\subsubsection*{22.}

If we have three integrals $a$, $b$, $c$, then the differential system may be written
\[
\frac{dx}{\,\partial(a, b, c)/\partial(y, z, w)\,} = \cdots = \frac{dw}{\,\partial(a, b, c)/\partial(x, y, z)\,},
\]
and the condition for a fourth integral is that the expression
\[
\frac{\partial(a, b, c)}{\partial(y, z, w)}\,dx + \cdots + \frac{\partial(a, b, c)}{\partial(x, y, z)}\,dw
\]
should be an exact differential; or, what is the same thing, that the Pfaffian derivatives of this expression should all vanish.

\subsection*{Particular Cases. Art.~Nos.~23 to~29.}

\subsubsection*{23.}

In the case of three variables $(x, y, z)$, the differential system is
\[
\frac{dx}{X} = \frac{dy}{Y} = \frac{dz}{Z},
\]
and the partial differential equation is
\[
X\frac{df}{dx} + Y\frac{df}{dy} + Z\frac{df}{dz} = 0.
\]

\newpage
\setcounter{page}{123}

\subsubsection*{24.}

In the case of three variables $(x, y, z)$, if we have one integral $a = \text{const.}$, then the differential system may be written
\[
\frac{dx}{\,\partial a/\partial y\,} = \frac{dy}{\,\partial a/\partial z\,} = \frac{dz}{\,\partial a/\partial x\,},
\]
and the condition for a second integral is that the expression
\[
\frac{\partial a}{\partial y}\,dx + \frac{\partial a}{\partial z}\,dy + \frac{\partial a}{\partial x}\,dz
\]
should be an exact differential.

\subsubsection*{25.}

If we have two integrals $a$, $b$ in the case of three variables, then
\[
\frac{dx}{\,\partial(a, b)/\partial(y, z)\,} = \frac{dy}{\,\partial(a, b)/\partial(z, x)\,} = \frac{dz}{\,\partial(a, b)/\partial(x, y)\,},
\]
and the condition for a third integral is that
\[
\frac{\partial(a, b)}{\partial(y, z)}\,dx + \frac{\partial(a, b)}{\partial(z, x)}\,dy + \frac{\partial(a, b)}{\partial(x, y)}\,dz
\]
should be an exact differential.

\subsubsection*{26.}

If the expression $X\,dx + Y\,dy + Z\,dz$ is an exact differential, say $= du$, then $u = \text{const.}$ is an integral; and the condition for this is
\[
\frac{dX}{dy} = \frac{dY}{dx}, \quad \frac{dY}{dz} = \frac{dZ}{dy}, \quad \frac{dZ}{dx} = \frac{dX}{dz}.
\]

\subsubsection*{27.}

The method of multipliers for three variables is as follows: if $\lambda$, $\mu$, $\nu$ be factors such that
\[
\lambda\,dx + \mu\,dy + \nu\,dz = 0,
\]
then $\lambda X + \mu Y + \nu Z = 0$, and the condition that this should be an exact differential gives rise to the equations for the determination of $\lambda$, $\mu$, $\nu$.

\newpage
\setcounter{page}{124}

\subsubsection*{28.}

If we have one integral $a = \text{const.}$ in the case of three variables, then the differential equation may be written
\[
\frac{dx}{\partial a/\partial y} = \frac{dy}{\partial a/\partial z} = \frac{dz}{\partial a/\partial x},
\]
and the condition for a second integral is that the expression
\[
\frac{\partial a}{\partial y}\,dx + \frac{\partial a}{\partial z}\,dy + \frac{\partial a}{\partial x}\,dz
\]
should be an exact differential; viz.~that we should have
\[
\frac{d}{dz}\biggl(\frac{\partial a}{\partial y}\biggr) = \frac{d}{dy}\biggl(\frac{\partial a}{\partial x}\biggr), \quad \&\text{c.}
\]

\subsubsection*{29.}

And we thence find
\[
\frac{d(a, b)}{d(x, y)} = \frac{d(a, b)}{d(y, z)} = \frac{d(a, b)}{d(z, x)} = 0,
\]
and each of these, as containing the differential coefficients of $a$ and $b$, gives the condition for the dependency of these functions.

\subsection*{Pfaffian Theorem. Art.~No.~30.}

\subsubsection*{30.}

According as the variables are odd or even in number, the expression
\[
X\,dx + Y\,dy + Z\,dz + W\,dw + \cdots
\]
may be of a more special form; for instance, it may be a sum of two terms $= \xi\,d\rho + \eta\,d\sigma$, or, what is the same thing, it may be reducible to the form
\[
X_{1}\,dx_{1} + Y_{1}\,dy_{1} + Z_{1}\,dz_{1} + W_{1}\,dw_{1} + \cdots,
\]
where the number of terms is less than the original number.

\newpage
\setcounter{page}{125}

The general theorem is that any such expression is reducible to a form in which the number of terms is at most equal to the number of variables; and, if the variables are even in number, the expression may be reducible to a form in which the number of terms is less than the number of variables.

In particular, if the variables are four in number, the expression $X\,dx + Y\,dy + Z\,dz + W\,dw$ is reducible to the form $\xi\,d\rho + \eta\,d\sigma$, where $\xi$, $\eta$, $\rho$, $\sigma$ are functions of $(x, y, z, w)$.

The condition for this is that the Pfaffian derivative $1234$ should vanish; viz.~that
\[
12.34 + 13.42 + 14.23 = 0.
\]

And if the variables are six in number, the expression $X\,dx + Y\,dy + Z\,dz + W\,dw + U\,du + V\,dv$ is reducible to the form $\xi\,d\rho + \eta\,d\sigma + \zeta\,d\tau$, where $\xi$, $\eta$, $\zeta$, $\rho$, $\sigma$, $\tau$ are functions of the variables.

The condition for this is that the Pfaffian derivative $123456$ should vanish; viz.~that
\[
12.3456 + 13.4561 + 14.5612 + 15.6123 + 16.2345 = 0.
\]

\subsection*{The Hamiltonian System. Art.~Nos.~31 to~39.}

\subsubsection*{31.}

The Hamiltonian system for a bipartite set $(x, y, z;\; p, q, r)$ is
\[
\frac{dx}{dt} = \frac{dH}{dp}, \quad \frac{dy}{dt} = \frac{dH}{dq}, \quad \frac{dz}{dt} = \frac{dH}{dr}, \quad
\frac{dp}{dt} = -\frac{dH}{dx}, \quad \frac{dq}{dt} = -\frac{dH}{dy}, \quad \frac{dr}{dt} = -\frac{dH}{dz},
\]
where $H$ is a given function of the variables.

\subsubsection*{32.}

The system may be written in the form
\[
\frac{dx}{dH/dp} = \frac{dy}{dH/dq} = \frac{dz}{dH/dr} = \frac{dp}{-dH/dx} = \frac{dq}{-dH/dy} = \frac{dr}{-dH/dz} = dt,
\]
and we associate with this the linear partial differential equation
\[
\frac{dH}{dp}\frac{d\Theta}{dx} - \frac{dH}{dx}\frac{d\Theta}{dp} + \frac{dH}{dq}\frac{d\Theta}{dy} - \frac{dH}{dy}\frac{d\Theta}{dq} + \frac{dH}{dr}\frac{d\Theta}{dz} - \frac{dH}{dz}\frac{d\Theta}{dr} = 0.
\]

\newpage
\setcounter{page}{126}

\subsubsection*{33.}

It is to be observed that the Hamiltonian system is of the same form as the general differential system, if we write
\[
X = \frac{dH}{dp}, \quad Y = \frac{dH}{dq}, \quad Z = \frac{dH}{dr}, \quad
P = -\frac{dH}{dx}, \quad Q = -\frac{dH}{dy}, \quad R = -\frac{dH}{dz},
\]
and the partial differential equation is then
\[
X\frac{d\Theta}{dx} + Y\frac{d\Theta}{dy} + Z\frac{d\Theta}{dz} + P\frac{d\Theta}{dp} + Q\frac{d\Theta}{dq} + R\frac{d\Theta}{dr} = 0.
\]

\subsubsection*{34.}

The Hamiltonian system establishes between the variables a fivefold relation depending upon five arbitrary constants; in fact, regarding $(x, y, z, p, q, r)$ as the coordinates of a point in six-dimensional space, and starting from any given point, the differential equations determine the direction of passage to a consecutive point, and so on; the locus of the point is therefore a curve, or we have between the coordinates a fivefold relation involving five arbitrary constants.

\subsubsection*{35.}

The constants will be expressible as functions of the coordinates, viz.~the fivefold relation will be expressible in the form
\[
a = \text{const.}, \quad b = \text{const.}, \quad c = \text{const.}, \quad d = \text{const.}, \quad e = \text{const.},
\]
where $a$, $b$, $c$, $d$, $e$ denote respective functions of $(x, y, z, p, q, r)$; and the differential equations thus establish between the variables the fivefold relation $a = \text{const.}$, $b = \text{const.}$, $c = \text{const.}$, $d = \text{const.}$, $e = \text{const.}$

\subsubsection*{36.}

If in the Hamiltonian system we have two integrals $a$, $b$, these may be regarded as given functions of the variables; and we may then express the ratios $dx : dy : dz : dp : dq : dr$ by means of the Jacobians. In fact, regarding $a$, $b$ as given functions, we have
\[
\frac{dx}{\,d(a, b)/d(p, y, z, q, r)\,} = \cdots,
\]
where the notation is used to denote the series of functional determinants.

\newpage
\setcounter{page}{127}

may be of a more special form; for instance, it may be a sum of two terms $= \xi\,d\rho + \eta\,d\sigma$, or,

\subsubsection*{37.}

The general integral of the partial differential equation may be written in the form
\[
\Phi(a, b, c, d, e) = 0,
\]
where $\Phi$ is an arbitrary functional symbol; and this is equivalent to the system of five equations $a = \text{const.}$, \ldots, $e = \text{const.}$, involving as it should do five arbitrary constants.

\subsubsection*{38.}

If we have three integrals $a$, $b$, $c$ of the Hamiltonian system, then the differential system may be written
\[
\frac{dx}{\,d(a, b, c)/d(p, q, y, z, r)\,} = \cdots,
\]
and the condition for a fourth integral is that the expression
\[
\frac{d(a, b, c)}{d(p, q, y, z, r)}\,dx + \cdots
\]
should be an exact differential.

\subsubsection*{39.}

If we have four integrals $a$, $b$, $c$, $d$ of the Hamiltonian system, then the differential system may be written
\[
\frac{dx}{\,d(a, b, c, d)/d(p, q, r, y, z)\,} = \cdots,
\]
and the condition for a fifth integral is that the expression
\[
\frac{d(a, b, c, d)}{d(p, q, r, y, z)}\,dx + \cdots
\]
should be an exact differential.

\newpage
\setcounter{page}{128}

and if we herein attend to the terms which contain $dt$, we find that the expression
\[
dH - \frac{dH}{dp}\,dp - \frac{dH}{dq}\,dq - \frac{dH}{dr}\,dr + \frac{dH}{dx}\,dx + \frac{dH}{dy}\,dy + \frac{dH}{dz}\,dz
\]
vanishes identically; and the system may therefore be written in the form
\begin{align*}
dx &= \frac{dH}{dp}\,dt, \quad dy = \frac{dH}{dq}\,dt, \quad dz = \frac{dH}{dr}\,dt, \\
dp &= -\frac{dH}{dx}\,dt, \quad dq = -\frac{dH}{dy}\,dt, \quad dr = -\frac{dH}{dz}\,dt,
\end{align*}
which puts in evidence the integral $H = \text{const.}$

\subsection*{The Poisson-Jacobi Theorem peculiar to the Hamiltonian Form. Art.~Nos.~40 to~45.}

\subsubsection*{40.}

The theorem in question is that if $a$, $b$ are any two integrals of the Hamiltonian system, then the function $(a, b)$ is also an integral.

\subsubsection*{41.}

To prove this, we observe that if $a$ is an integral, then
\[
\frac{da}{dt} = \frac{dH}{dp}\frac{da}{dx} - \frac{dH}{dx}\frac{da}{dp} + \frac{dH}{dq}\frac{da}{dy} - \frac{dH}{dy}\frac{da}{dq} + \frac{dH}{dr}\frac{da}{dz} - \frac{dH}{dz}\frac{da}{dr} = 0,
\]
and similarly for $b$; and hence, differentiating the expression $(a, b)$, we find
\[
\frac{d}{dt}(a, b) = \biggl(\frac{da}{dt}, b\biggr) + \biggl(a, \frac{db}{dt}\biggr) = 0,
\]
which proves the theorem.

\newpage
\setcounter{page}{129}

\subsubsection*{42.}

The theorem may also be stated as follows: if $a$, $b$ are any two functions of the variables such that $(a, b) = 0$, and if one of them, say $a$, is an integral of the Hamiltonian system, then the other, $b$, is also an integral.

\subsubsection*{43.}

But we wish to prove the theorem directly from the definition of the symbol $(a, b)$; viz.~we wish to show that if $a$, $b$ are integrals, then $(a, b)$ is also an integral.

\subsubsection*{44.}

We have, by definition,
\[
(a, b) = \frac{d(a, b)}{d(p, x)} + \frac{d(a, b)}{d(q, y)} + \frac{d(a, b)}{d(r, z)},
\]
and we wish to show that if $a$, $b$ are integrals, then $(a, b)$ is also an integral.

\subsubsection*{45.}

To prove this, we observe that
\begin{align*}
\frac{d}{dt}(a, b) &= \sum\biggl\{\frac{d}{dt}\biggl(\frac{da}{dp}\frac{db}{dx} - \frac{da}{dx}\frac{db}{dp}\biggr)\biggr\} \\
&= \sum\biggl\{\biggl(\frac{d}{dt}\frac{da}{dp}\biggr)\frac{db}{dx} + \frac{da}{dp}\biggl(\frac{d}{dt}\frac{db}{dx}\biggr) - \biggl(\frac{d}{dt}\frac{da}{dx}\biggr)\frac{db}{dp} - \frac{da}{dx}\biggl(\frac{d}{dt}\frac{db}{dp}\biggr)\biggr\},
\end{align*}
where the summation extends to the three pairs $(p, x)$, $(q, y)$, $(r, z)$.

Now, since $a$ is an integral, we have
\[
\frac{da}{dt} = \frac{dH}{dp}\frac{da}{dx} - \frac{dH}{dx}\frac{da}{dp} + \cdots = 0,
\]
and hence, differentiating with respect to $p$, we find
\[
\frac{d}{dp}\frac{da}{dt} = \frac{d^{2}H}{dp^{2}}\frac{da}{dx} + \frac{dH}{dp}\frac{d^{2}a}{dp\,dx} - \frac{d^{2}H}{dp\,dx}\frac{da}{dp} - \frac{dH}{dx}\frac{d^{2}a}{dp^{2}} + \cdots = 0,
\]
and similarly for the other derivatives.

\newpage
\setcounter{page}{130}

Hence, unless $(a, b) = 0$, we have $A = \dfrac{1}{J}$. The remaining five equations give similarly $B = \dfrac{1}{J}$, $C = \dfrac{1}{J}$, $D = \dfrac{1}{J}$, $E = \dfrac{1}{J}$; and we have therefore
\[
(a, b) = \frac{d(a, b)}{d(p, x)} + \frac{d(a, b)}{d(q, y)} + \frac{d(a, b)}{d(r, z)} = \frac{3}{J},
\]
or, since the first two terms are $= \dfrac{2}{J}$, we have $(a, b) = \dfrac{1}{J}$.

\medskip\noindent\textsc{C.~X.}\hfill 14

\newpage
\setcounter{page}{131}

\subsection*{Further Development of the Hamiltonian System. Art.~Nos.~46 to~58.}

\subsubsection*{46.}

where $l$, $m$ are indeterminate functions of the variables; and we seek to determine these functions so that the expression
\[
l\,da + m\,db + n\,dc
\]
shall be an exact differential.

\subsubsection*{47.}

The condition for this is that the Pfaffian derivatives of the expression shall all vanish; viz.~that we shall have
\[
\frac{d(l, a)}{d(x, y)} + \frac{d(m, b)}{d(x, y)} + \frac{d(n, c)}{d(x, y)} = 0,
\]
and similarly for the other pairs of variables.

\subsubsection*{48.}

If we assume that $l$, $m$, $n$ are functions of $(a, b, c)$ only, then the conditions reduce to
\[
l\frac{d(b, c)}{d(x, y)} + m\frac{d(c, a)}{d(x, y)} + n\frac{d(a, b)}{d(x, y)} = 0,
\]
and similarly for the other pairs; and these conditions are satisfied if
\[
l : m : n = \frac{d(b, c)}{d(x, y, z)} : \frac{d(c, a)}{d(x, y, z)} : \frac{d(a, b)}{d(x, y, z)}.
\]

\newpage
\setcounter{page}{132}

\subsubsection*{49.}

But we may also consider the case where $l$, $m$ are functions of the variables, not necessarily of $(a, b, c)$ only; and in this case the conditions are more complicated.

\subsubsection*{50.}

The general theory of the Hamiltonian system may be developed as follows: if $a$, $b$, $c$, $d$, $e$ are any five integrals, then the system may be written
\[
\frac{dx}{\,d(a, b, c, d, e)/d(p, q, r, y, z)\,} = \cdots,
\]
and the condition for a sixth integral is that the expression
\[
\frac{d(a, b, c, d, e)}{d(p, q, r, y, z)}\,dx + \cdots
\]
should be an exact differential.

\newpage
\setcounter{page}{133}

Hence, unless $d(a, b)/d(x, y) = 0$, we have $A = \dfrac{1}{J}$. The remaining five equations give similarly
\[
B = \frac{1}{J}, \quad C = \frac{1}{J}, \quad D = \frac{1}{J}, \quad E = \frac{1}{J};
\]
and we have therefore
\[
(a, b) = \frac{d(a, b)}{d(p, x)} + \frac{d(a, b)}{d(q, y)} + \frac{d(a, b)}{d(r, z)} = \frac{3}{J},
\]
or, since the first two terms are $= \dfrac{2}{J}$, we have $(a, b) = \dfrac{1}{J}$.

\subsubsection*{51.}

The theorem may be generalized as follows: if $a$, $b$, $c$ are any three integrals of the Hamiltonian system, then the function
\[
\frac{d(a, b, c)}{d(p, q, x)} + \frac{d(a, b, c)}{d(p, r, x)} + \cdots
\]
is also an integral.

\subsubsection*{52.}

And more generally, if $a$, $b$, $c$, \ldots\ are any number of integrals, then the function formed by summing the functional determinants
\[
\frac{d(a, b, c, \ldots)}{d(p, q, r, \ldots, x, y, z, \ldots)}
\]
is also an integral.

\newpage
\setcounter{page}{134}

or, since the first two terms are equal, we have
\[
(a, b) = \frac{d(a, b)}{d(p, x)} + \frac{d(a, b)}{d(q, y)} + \frac{d(a, b)}{d(r, z)} = 0.
\]

\subsubsection*{53.}

The theorem may be stated in a different form as follows: if $a$, $b$ are any two integrals of the Hamiltonian system, then the expression
\[
\frac{d(a, b)}{d(p, x)} + \frac{d(a, b)}{d(q, y)} + \frac{d(a, b)}{d(r, z)}
\]
is also an integral; and conversely, if this expression vanishes, then $a$, $b$ are integrals.

\subsubsection*{54.}

The condition $(a, b) = 0$ is the condition that the integrals $a$, $b$ should be ``in involution''; and the theorem is that if two integrals are in involution, then any third integral is also in involution with each of them.

\newpage
\setcounter{page}{135}

\subsubsection*{50.}

Supposing that any two integrals $a$, $b$ of the Hamiltonian system are given, we can find a third integral $c$ such that $(a, c) = 0$ and $(b, c) = 0$; and we can then find a fourth integral $d$ such that $(a, d) = 0$, $(b, d) = 0$, $(c, d) = 0$; and so on.

\subsubsection*{51.}

The process may be continued until we have obtained a complete system of integrals, viz.~a system of $n$ integrals such that any two of them are in involution.

\subsubsection*{52.}

If we have a complete system of $n$ integrals, then the Hamiltonian system may be reduced to the form
\[
\frac{dx_{1}}{dt} = \frac{dH}{dp_{1}}, \quad \frac{dy_{1}}{dt} = \frac{dH}{dq_{1}}, \quad \ldots,
\]
where $x_{1}$, $y_{1}$, \ldots, $p_{1}$, $q_{1}$, \ldots\ are functions of the original variables.

\newpage
\setcounter{page}{136}

\subsubsection*{53.}

We consider the differential expression
\[
dV - p\,dx - q\,dy - r\,dz,
\]
which, treating $V$ as a function of $(x, y, z)$, and $p$, $q$, $r$ as differential coefficients, may be written
\[
dV = p\,dx + q\,dy + r\,dz.
\]

\subsubsection*{54.}

If $V$ is a function of $(x, y, z)$, then $p = \dfrac{dV}{dx}$, $q = \dfrac{dV}{dy}$, $r = \dfrac{dV}{dz}$; and the expression $dV - p\,dx - q\,dy - r\,dz$ vanishes identically.

\subsubsection*{55.}

But if $V$ is not a function of $(x, y, z)$ alone, but contains also other variables, then the expression does not vanish identically; and we may consider it as a function of the variables.

\newpage
\setcounter{page}{137}

and hence $A = (t - T)^{2} + A'$, and the like for $B$, $C$, \&c.

\subsubsection*{56.}

If we have $V = V(x, y, z, a, b)$, where $a$, $b$ are constants, then
\[
dV = \frac{dV}{dx}\,dx + \frac{dV}{dy}\,dy + \frac{dV}{dz}\,dz + \frac{dV}{da}\,da + \frac{dV}{db}\,db,
\]
and if $p = \dfrac{dV}{dx}$, $q = \dfrac{dV}{dy}$, $r = \dfrac{dV}{dz}$, we have
\[
dV - p\,dx - q\,dy - r\,dz = \frac{dV}{da}\,da + \frac{dV}{db}\,db.
\]

\subsubsection*{57.}

If $a$, $b$ are not constants but variables, then the expression on the right-hand side is not zero; and we may consider the equation
\[
dV - p\,dx - q\,dy - r\,dz = 0
\]
as a differential equation determining $V$ as a function of $(x, y, z)$.

\newpage
\setcounter{page}{138}

and again treating $a$, $b$, $H$ as variable, we have
\[
dV - p\,dx - q\,dy - r\,dz = dK + (t - T)\,dH + A\,da + B\,db,
\]
where
\begin{align*}
K &= V - px - qy - rz + T H, \\
t &= \frac{dV}{dH}, \quad T = \frac{dK}{dH}, \\
A &= \frac{dV}{da} - p\frac{dx}{da} - q\frac{dy}{da} - r\frac{dz}{da} + T\frac{dH}{da}, \\
B &= \frac{dV}{db} - p\frac{dx}{db} - q\frac{dy}{db} - r\frac{dz}{db} + T\frac{dH}{db}.
\end{align*}

\subsection*{The Partial Differential Equation $H = \text{const.}$ Art.~Nos.~59 to~70.}

\subsubsection*{59.}

We consider the partial differential equation
\[
H(x, y, z, p, q, r) = \text{const.},
\]
where $p$, $q$, $r$ are the differential coefficients of a function $V$ of $(x, y, z)$; viz.~$p = \dfrac{dV}{dx}$, $q = \dfrac{dV}{dy}$, $r = \dfrac{dV}{dz}$.

\newpage
\setcounter{page}{139}

\subsubsection*{60.}

The equation $H = \text{const.}$ is equivalent to the system of ordinary differential equations
\[
\frac{dx}{dH/dp} = \frac{dy}{dH/dq} = \frac{dz}{dH/dr} = \frac{dp}{-dH/dx} = \frac{dq}{-dH/dy} = \frac{dr}{-dH/dz},
\]
which is the Hamiltonian system.

\subsubsection*{61.}

Conversely, if we have the Hamiltonian system, and if we can find a function $V$ such that $p = \dfrac{dV}{dx}$, $q = \dfrac{dV}{dy}$, $r = \dfrac{dV}{dz}$, then $H = \text{const.}$ is a partial differential equation satisfied by~$V$.

\subsubsection*{62.}

The Hamiltonian system determines the variables as functions of $t$ and of five arbitrary constants; and if we can express these in terms of $x$, $y$, $z$, and substitute in the expression $p\,dx + q\,dy + r\,dz$, we obtain the differential of the function~$V$.

\newpage
\setcounter{page}{140}

Adding these together, we have
\[
\frac{d(A, B)}{d(x, y)} + \frac{d(q, r)}{d(y, z)} + \frac{d(r, p)}{d(z, x)} + \frac{d(p, q)}{d(x, y)} = 0,
\]
and, if it vanish, it is a condition that the expression $p\,dx + q\,dy + r\,dz$ should be an exact differential.

\subsubsection*{63.}

The condition may be written in the form
\[
(p, q) + (q, r) + (r, p) = 0,
\]
where $(p, q)$ denotes the functional determinant $\dfrac{d(p, q)}{d(x, y)}$, \&c.

\newpage
\setcounter{page}{141}

and, if it vanish, it is

\subsubsection*{64.}

The general solution of the partial differential equation $H = \text{const.}$ may be obtained as follows: we first integrate the Hamiltonian system, obtaining the variables as functions of $t$ and five arbitrary constants; we then express $p$, $q$, $r$ in terms of $x$, $y$, $z$, and substitute in the expression $p\,dx + q\,dy + r\,dz$; if this is an exact differential, its integral gives the function~$V$.

\subsubsection*{65.}

The process may be simplified if we know any integral of the Hamiltonian system; for if $a$ is such an integral, we may use it to reduce the order of the system.

\newpage
\setcounter{page}{142}

a solution of the equation $d\psi = \text{const.}$

\subsubsection*{66.}

If we have a complete system of $n$ integrals of the Hamiltonian system, then the partial differential equation $H = \text{const.}$ may be solved by quadratures.

\subsubsection*{67.}

In particular, if $n = 3$, and if we have three integrals $a$, $b$, $c$ in involution, then the solution may be obtained as follows: we form the expression
\[
\Theta = \int(p\,dx + q\,dy + r\,dz),
\]
and determine the arbitrary functions which appear in the integration by the condition that $\Theta$ shall be a function of $(x, y, z)$ only.

\newpage
\setcounter{page}{143}

we find for $\psi$ the equation
\[
\frac{d\psi}{dp} = \frac{d\psi}{dq} = \frac{d\psi}{dr} = 0,
\]
and hence $\psi$ is a function of $(x, y, z)$ only.

\subsubsection*{68.}

The function $\psi$ being determined, the general solution of the partial differential equation is
\[
V = \Theta + \psi,
\]
where $\Theta$ is the particular integral and $\psi$ is the arbitrary function.

\newpage
\setcounter{page}{144}

and each of these, as containing

\subsubsection*{69.}

The general integral of the partial differential equation $H = \text{const.}$ may be expressed in the form
\[
V = \int(p\,dx + q\,dy + r\,dz) + \Phi(a, b, c),
\]
where $\Phi$ is an arbitrary functional symbol, and $a$, $b$, $c$ are integrals of the Hamiltonian system.

\subsubsection*{70.}

The solution thus involves three arbitrary functions, which may be determined by giving the values of $V$, $\dfrac{dV}{dx}$, $\dfrac{dV}{dy}$ at any point.

\newpage
\setcounter{page}{145}

which, in virtue of $(a, H) = 0$, $(b, H) = 0$, becomes
\[
\frac{d(a, b)}{d(p, x)} + \frac{d(a, b)}{d(q, y)} + \frac{d(a, b)}{d(r, z)} = 0,
\]
and this is the condition that $a$, $b$ should be in involution.

\subsection*{Further Theory of the Partial Differential Equation $H = \text{const.}$ Art.~Nos.~71 to~80.}

\subsubsection*{71.}

The partial differential equation $H = \text{const.}$ may be treated by the method of characteristics; the characteristics are the curves determined by the Hamiltonian system.

\newpage
\setcounter{page}{146}

We have as a conjugate system

\subsubsection*{72.}

If $a$, $b$ are any two integrals of the Hamiltonian system which are in involution, then the equations $a = \text{const.}$, $b = \text{const.}$ determine a family of surfaces which are characteristics of the partial differential equation.

\subsubsection*{73.}

The characteristics may be used to construct the general solution of the partial differential equation; for if we know the values of $V$ and its differential coefficients on any surface which is not a characteristic, we can determine the solution in the neighbourhood of this surface.

\newpage
\setcounter{page}{147}

where the term in $d\theta$ is

which, from the equation which determines $\theta$,

\subsubsection*{74.}

The method of characteristics may be extended to the case where the variables are any even number $2n$; the Hamiltonian system then determines the characteristics, and the general solution may be obtained by quadratures.

\subsubsection*{75.}

If we have a complete system of $n$ integrals in involution, then the partial differential equation may be solved by quadratures; and the solution involves $n$ arbitrary functions.

\newpage
\setcounter{page}{148}

Forming from these the expression for $dV - p\,dx - q\,dy - r\,dz$,

and again treating $a$, $b$, $H$ as variable, we have
\[
dV - p\,dx - q\,dy - r\,dz = dK + (t - T)\,dH + A\,da + B\,db,
\]
where
\begin{align*}
K &= V - px - qy - rz + T H, \\
t &= \frac{dV}{dH}, \quad T = \frac{dK}{dH}, \\
A &= \frac{dV}{da} - p\frac{dx}{da} - q\frac{dy}{da} - r\frac{dz}{da} + T\frac{dH}{da}, \\
B &= \frac{dV}{db} - p\frac{dx}{db} - q\frac{dy}{db} - r\frac{dz}{db} + T\frac{dH}{db}.
\end{align*}

\newpage
\setcounter{page}{149}

expressions $x - p$, $y - q$, $z - r$ is an integral, or again the product of any expression of the first set into any expression of the second set is an integral.

\subsubsection*{76.}

The general theory may be applied to the case where $H$ is a homogeneous function of $(p, q, r)$; in this case the Hamiltonian system may be simplified by the introduction of new variables.

\subsubsection*{77.}

If $H$ is homogeneous of degree $k$ in $(p, q, r)$, then we have
\[
p\frac{dH}{dp} + q\frac{dH}{dq} + r\frac{dH}{dr} = kH,
\]
and the Hamiltonian system may be written in a form which puts this homogeneity in evidence.

\newpage
\setcounter{page}{150}

whence
\begin{multline}
dV - p\,dx - q\,dy - r\,dz = d\lambda + (t - T)\,dH + a'\,da + b'\,db + c'\,dc \\
- a\,da' - b\,db' - c\,dc';
\end{multline}
or, in place of

\medskip\noindent\textsc{C.~X.}\hfill 15

\end{document}
